
\section{The Area of Work}
Large Language Models (LLMs) are now used in chatbots, coding assistants, customer support systems and
personal assistants that people talk to over days, weeks and months. The underlying Transformer
architecture, however, is stateless: every request is answered only from the tokens that fit in the current
context window. Everything that was said in an earlier session is forgotten unless it is sent again.

This project lies at the intersection of Large Language Models, text embeddings, vector databases,
information retrieval and knowledge graphs. The first phase of the project (April 2026) analysed the
limitations of the context window and of Retrieval-Augmented Generation (RAG), and \emph{proposed} a
persistent memory layer. This second phase \emph{builds} that memory layer, benchmarks it against the
alternatives, and reports what works and what does not.

\section{Problem Addressed}
In Phase 1 we established, using the RAGTruth dataset and Needle-in-a-Haystack (NIAH) experiments, two
limitations of current systems:
\begin{itemize}[nosep]
  \item LLMs do not use long contexts effectively. Accuracy drops as the context grows and depends on
        \emph{where} the relevant information sits (``lost in the middle'').
  \item Retrieval alone does not remove hallucination. Hallucination detectors on RAGTruth reached high
        recall (0.88--0.95) but very low precision (0.25--0.26), i.e.\ an F1 of only about 0.4.
\end{itemize}
Phase 1 therefore proposed a memory layer that stores conversation turns as embeddings in a vector
database and injects only the most relevant memories into the prompt. Its future-work section listed the
concrete components that were still missing: hybrid vector + BM25 search, a multi-stage retrieval pipeline
with re-ranking, query expansion, memory scoring with ageing, a graph-based memory representation, an
end-to-end pipeline, and a comparative evaluation against baseline LLMs and RAG.

The problem addressed in this phase is to turn that proposal into a working, measurable system and to
answer, with experiments, whether a memory layer actually gives better grounding, better efficiency and
better long-term consistency than (a) a plain LLM, (b) a larger context window and (c) plain vector RAG.

\section{Motivation}
A chat assistant that re-sends the whole history on every turn pays for thousands of tokens per query and,
as Phase 1 showed, still misses facts buried in the middle of that history. A sliding window of recent turns
is cheap but forgets anything older than a few sessions. Vector RAG over the history is cheap and retrieves
well when the question and the memory share wording, but it has no notion of \emph{which} facts are current,
\emph{how} facts are connected through people and entities, or \emph{whether} the generated answer is
actually supported by what was retrieved.

A memory layer that selects a handful of relevant, current, connected notes and verifies the answer against
them promises to be simultaneously cheaper, more accurate and less prone to hallucination. That promise
needs to be tested rather than assumed, which is what this phase does.

\section{Objectives}
\begin{itemize}[nosep]
  \item Implement the complete memory layer proposed in Phase 1 as a modular Python package: embedding
        store (FAISS), BM25 keyword index, reciprocal-rank fusion, cross-encoder re-ranking, entity graph,
        memory scoring, conflict-aware updates and token-budgeted context injection.
  \item Add a grounding gate that checks whether a generated answer is supported by the injected memories,
        and compare it with the Phase-1 hallucination-detection numbers.
  \item Build a long-term conversational memory benchmark with single-hop, temporal (updated fact),
        multi-hop and unanswerable questions, and a NIAH-style context-length benchmark.
  \item Compare, with the same local LLM, five answering strategies: no memory, full conversation history,
        sliding window, plain vector RAG (the Phase-1 design) and the proposed memory layer.
  \item Quantify the contribution of each pipeline stage through an ablation study and study the trade-off
        between memory pruning and recall.
\end{itemize}

\section{Existing System}
\subsection{Retrieval-Augmented Generation (RAG)}
RAG converts documents or conversation chunks into embeddings, stores them in a vector database and, at
query time, prepends the top-$k$ nearest chunks to the prompt. It is the Phase-1 baseline design. RAG is
stateless with respect to time: it does not know that ``I moved to Pune'' supersedes ``I live in Jaipur'',
and it retrieves each chunk independently, so facts that are only connected through a person's name are
easily missed.

\subsection{Context Window Expansion}
Modern models accept 32k to over a million tokens. Sending the whole conversation history is the simplest
form of memory, but Phase 1 showed that effective context is much shorter than the advertised window, cost
grows linearly with history length, and accuracy decays with length and position. We keep this approach as
the ``full context'' baseline in all Phase-2 experiments.

\subsection{Limitations of Existing Systems}
Neither approach provides persistence with \emph{selectivity}: RAG is selective but not time- or
entity-aware, while the full context is complete but neither selective nor efficient. Neither verifies that
the final answer is grounded. The memory layer built in this phase is designed to close exactly these gaps.

\section{Organisation of the Report}
Chapter 2 reviews the literature, including the memory-system work that became relevant in this phase.
Chapter 3 describes the design and implementation of the memory layer and the benchmarks. Chapter 4 reports
the simulation results, and Chapter 5 concludes and lists future work.
